% Gesamttest «Zentrische Streckung und Ähnlichkeit» — Quelle des PDFs (python3 scripts/build-lp-pdf.py aehnlichkeit).
% Musterlösung und Raster: bewertungspaket.tex. Alle Werte mit python3 nachgerechnet (08.10.2026, scripts/lp/aehnlichkeit/zahlen.py).
% Kein \lt/\gt — pdfLaTeX kennt sie nicht, < und > tun es.
\documentclass[11pt]{article}
\usepackage{../lp-druck}
\renewcommand{\lpname}{Zentrische Streckung und Ähnlichkeit}
\renewcommand{\lpteilgebiet}{5.2}
\renewcommand{\lpbereich}{Grundlagenfach}
\renewcommand{\lpfuss}{Gesamttest · 25 Punkte}
\begin{document}
\lpkopf{Gesamttest}{%
  \vspace{4pt}\noindent\begin{tabularx}{\linewidth}{@{}X X@{}}
  Code (kein Name): \hrulefill & Punkte: \hrulefill\ / 25
  \end{tabularx}}

\begin{lpinfo}{Anleitung}
\textbf{Zeit:} rund 30 Minuten. \textbf{Hilfsmittel:} Taschenrechner erlaubt, für alle Aufgaben (RLP 5.2 trägt keinen Vermerk
«auch ohne Hilfsmittel»). Schreib den \textbf{Rechenweg} auf — bewertet wird der Weg; zu Sachaufgaben gehört eine
\textbf{Skizze}. Gerundete Ergebnisse auf zwei Dezimalen, Zwischenresultate ungerundet weiterverwenden.
Figuren sind Skizzen, nicht massstäblich.
Danach: Lösung scannen oder fotografieren (ohne Namen und Standort) und mit dem \emph{Bewertungspaket} bewerten lassen.
\end{lpinfo}

\begin{lpaufgabe}{G1}{Zentrische Streckung mit Koordinaten}{4 P}
Das Dreieck $P(4 \mid 0)$, $Q(6 \mid 3)$, $R(3 \mid 2)$ wird vom Zentrum $Z(2 \mid -1)$ aus mit $k = -1.5$ gestreckt.
\begin{enumerate}[label=(\alph*),leftmargin=*,itemsep=2pt]
\item Berechne die Bildpunkte $P'$ und $Q'$.
\item Die Seite $PQ$ ist $\sqrt{13} \approx 3.61\,\text{cm}$ lang. Wie lang ist $P'Q'$?
\item Welche Seite des Bilddreiecks ist parallel zu $QR$? Begründe in einem Satz.
\end{enumerate}
\lpplatz{6}
\end{lpaufgabe}

\begin{lpaufgabe}{G2}{Strahlensätze}{4 P}
\begin{minipage}[t]{0.52\linewidth}\vspace{0pt}
Im Bild ist $AB \parallel A'B'$. Es ist $\overline{SA} = 2.5\,\text{cm}$, $\overline{AA'} = 3.5\,\text{cm}$,
$\overline{SB} = 3\,\text{cm}$ und $\overline{AB} = 2\,\text{cm}$.
\begin{enumerate}[label=(\alph*),leftmargin=*,itemsep=2pt]
\item Berechne $\overline{BB'}$.
\item Berechne $\overline{A'B'}$.
\end{enumerate}
\end{minipage}\hfill
\begin{minipage}[t]{0.44\linewidth}\vspace{0pt}\centering
\begin{tikzpicture}[scale=0.62]
  \draw[lpgrau] (-0.4,0) -- (7,0); \draw[lpgrau] (-0.3,-0.265) -- (6,5.29);
  \draw[lpblau,thick] (2.5,0) -- (2.25,1.984); \draw[lporange,thick] (6,0) -- (5.4,4.762);
  \foreach \p in {(0,0),(2.5,0),(6,0),(2.25,1.984),(5.4,4.762)} \fill \p circle (2pt);
  \node[below] at (0,0) {$S$}; \node[below] at (2.5,0) {$A$}; \node[below] at (6,0) {$A'$};
  \node[above left] at (2.25,1.984) {$B$}; \node[above left] at (5.4,4.762) {$B'$};
\end{tikzpicture}
\end{minipage}
\lpplatz{5}
\end{lpaufgabe}

\begin{lpaufgabe}{G3}{Lochkamera}{3 P}
Eine Lochkamera ist eine dunkle Schachtel mit einem kleinen Loch in der Vorderwand. Auf der Rückwand, $15\,\text{cm}$ hinter
dem Loch, entsteht ein umgekehrtes Bild. Ein $1.8\,\text{m}$ grosser Mensch steht $6\,\text{m}$ vor dem Loch.
\begin{enumerate}[label=(\alph*),leftmargin=*,itemsep=2pt]
\item Skizziere die Lichtstrahlen vom Kopf und von den Füssen durch das Loch. Warum darfst du die Strahlensätze anwenden?
\item Wie gross ist das Bild des Menschen?
\item Wie weit muss der Mensch vom Loch entfernt stehen, damit sein Bild $3\,\text{cm}$ gross ist?
\end{enumerate}
\lpplatz{7}
\end{lpaufgabe}

\begin{lpaufgabe}{G4}{Karte}{4 P}
Eine Wanderkarte hat den Massstab $1 : 25\,000$.
\begin{enumerate}[label=(\alph*),leftmargin=*,itemsep=2pt]
\item Ein Wanderweg ist auf der Karte $18.4\,\text{cm}$ lang. Wie lang ist er in Wirklichkeit (in km)?
\item Ein See bedeckt auf der Karte $3.2\,\text{cm}^2$. Wie gross ist er in Wirklichkeit (in $\text{km}^2$)?
\item Wie viele $\text{cm}^2$ bedeckt derselbe See auf einer Karte im Massstab $1 : 50\,000$?
\end{enumerate}
\lpplatz{6}
\end{lpaufgabe}

\begin{lpaufgabe}{G5}{Blumenbeet}{3 P}
Ein kreisrundes Blumenbeet hat $1.6\,\text{m}$ Durchmesser. Es soll durch ein kreisrundes Beet mit der dreifachen Fläche
ersetzt werden.
\begin{enumerate}[label=(\alph*),leftmargin=*,itemsep=2pt]
\item Welchen Durchmesser hat das neue Beet?
\item Mit welchem Faktor wächst der Umfang (die Länge der Einfassung)?
\end{enumerate}
\lpplatz{5}
\end{lpaufgabe}

\begin{lpaufgabe}{G6}{Zwei Dreiecke}{4 P}
Im Dreieck $ABC$ ist $\alpha = 48^\circ$ und $\beta = 64^\circ$; die Seiten sind $\overline{AB} = 6\,\text{cm}$,
$\overline{BC} \approx 4.81\,\text{cm}$ und $\overline{CA} \approx 5.82\,\text{cm}$. Im Dreieck $PQR$ ist der Winkel bei $P$
$68^\circ$, der Winkel bei $Q$ $48^\circ$, und $\overline{QR} = 9\,\text{cm}$.
\begin{enumerate}[label=(\alph*),leftmargin=*,itemsep=2pt]
\item Begründe, warum die Dreiecke ähnlich sind, und gib an, welche Ecken einander entsprechen.
\item Wie lang ist $\overline{PQ}$?
\item Wie viel Mal so gross wie die Fläche von $ABC$ ist die Fläche von $PQR$?
\end{enumerate}
\lpplatz{7}
\end{lpaufgabe}

\begin{lpaufgabe}{G7}{Rechtwinkliges Dreieck}{3 P}
Im Dreieck $ABC$ liegt der rechte Winkel bei $C$. Die Hypotenuse ist $c = 12.5\,\text{cm}$ lang, die Kathete
$a = \overline{BC} = 7.5\,\text{cm}$. Die Höhe $h$ auf $c$ teilt die Hypotenuse in $p$ (an $a$) und $q$ (an $b$).
\begin{enumerate}[label=(\alph*),leftmargin=*,itemsep=2pt]
\item Berechne $p$ und $q$.
\item Berechne die Höhe $h$.
\end{enumerate}
\lpplatz{5}
\end{lpaufgabe}
\end{document}
